\cotefile{34}
\scanpage{50}
\mspage{48}
\margin{(\unsure{Deligne})}
% lines 1-7
\inserted{Proposition} \deleted{\uline{Corollaire}} 9.6. Supposons résolution des singularités sur $K^{\infty}$ en ttes dimensions, et que pour tt $X$ \deleted{\unsure{qui}} \uline{projectif et lisse} sur $K$, \deleted{\unsure{les}} l'action de $I$ sur $\unsure{H^i_q|H^i_c}(X_{\overline K},\mathbb{Q}_\ell)$ soit quasi-unipotente, et \unsure{kif-kif} après extension finie radicielle de $K$. Alors \deleted{\unsure{l'action de}} pour tt $X_K$ de t.f. sur $K$, l'action de $I$ sur les $H^i(X_{\overline K},\mathbb{Q}_\ell)$ est quasi-unipotente.

% lines 8-9
Il suffit de \unsure{conjuguer|combiner} \deleted{\unsure{9.6}} \inserted{\unsure{ce|la} qui précède} et \unsure{9.4. c)}, en procédant par récurrence sur \unsure{$n$}.

% lines 10-14
\deleted{\uline{Corollaire} \unsure{9.6.}}

9.7. \uline{Cas des} $H^i(X_{\overline K},\mathbb{Q}_\ell)$, pour $i\leq 1$.

Le cas de $i=0$ est trivial et laissé au lecteur (en fait, $\pi$ y opère de façon \uuline{quasi-triviale}).

% lines 15-21
Pour le cas $i=1$, on \deleted{\unsure{se ramène}} peut, quitte à passer à une extension finie de $K$, supposer que les composantes irréductibles réduites de $X_K$ sont géom. intègres. Supposons $X$ réduit, soit $X'$ sa normalisée, $X''=X'\times_X X'$, \inserted{$X'''=X'\times_X X'\times_X X'$}, \unsure{etc.} Alors, on \deleted{\unsure{a}} a une suite spectrale de descente finie (SGA 4 \uline{VIII}) :
% lines 22-24
\[
\deleted{0\longrightarrow\operatorname{Ker}\bigl(H^0(X''_{\overline K})\overset{\partial}{\longrightarrow}H^0(X'''_{\overline K})\bigr)}
\]
\[
\begin{aligned}
0\longrightarrow H^1\bigl(\alpha\longmapsto H^0((\unsure{X'|X^\bullet}/X)^{\alpha+1}_{\overline K},\mathbb{Q}_\ell)\bigr)&\longrightarrow H^1(X_{\overline K},\mathbb{Q}_\ell)\deleted{\unsure{\longrightarrow H^1}}\\
&\longrightarrow\operatorname{Ker}\bigl(H^1(X'_{\overline K},\mathbb{Q}_\ell)\overset{\partial}{\underset{\unsure{\neq|\delta}}{\longrightarrow}}H^1(X''_{\overline K},\mathbb{Q}_\ell)\bigr)
\end{aligned}
\]
\margin{\unsure{quotient d'un sous-espace} de $H^0(X''_{\overline K},\mathbb{Q}_\ell)$}
% lines 25-26
Cela nous ramène donc à voir que l'action de $I$ sur \deleted{\unsure{$X'$}} $H^1(X'_{\overline K},\mathbb{Q}_\ell)$. Donc on est ramené au
% archive sheet number: 49
\scanpage{51}
\mspage{49}
% lines 1-4
cas de $X$ géom. normal. \deleted{\unsure{Le résultat est} \unsure{$U\subset X$}} \inserted{On peut supposer de} plus $X$ géom. intègre (\deleted{\unsure{en raisonnant}} les composantes connexes l'étant). Si $U$ est un ouvert non vide de $X$, on a $H^1(X_{\overline K},\mathbb{Q}_\ell)\hookrightarrow H^1(U_{\overline K},\mathbb{Q}_\ell)$, donc on \deleted{peut}

% lines 5-13: large crossed block
\deleted{supposer $X$ affine, donc quasi-projectif. Soit $\hat X$ \unsure{est associé à} $U$. \unsure{Pour} $X$, \unsure{Et on peut donc} supposer $X$ affine, donc quasi-projectif. Quitte à faire encore une extension \inserted{finie} radicielle sur $K$, on peut supposer que $X\subset\overline X$, où $\overline X$ est projectif, géom. normal. Notons que \unsure{pour remplacer} $X$ par $U$. Prenant $U$ assez petit, $U$ est un revêtement étale fini d'un ouvert $V$ de $\mathbb{P}^n$. Alors on a $H^1(\unsure{\cdot}$}

% lines 14-20: further cancelled attempt, with additions
\deleted{\unsure{remplacer}} peut \deleted{\unsure{supposer} $X$ affine par $U$, qu'on peut supposer} \inserted{$X$} affine \inserted{et lisse}.

\deleted{Soit \unsure{Considérons} une clôture projective de \inserted{$X$} $U$, qu'on peut supposer (quitte à faire encore une extension radicielle $K'/K$, finie) géom. normale. Alors $H^1(U,\mathbb{Z}/\ell^\nu\mathbb{Z})$ classifie les \unsure{rev. galoisiens} \illegible{2} de $X$, \unsure{mais} \unsure{normaux} sur $X$, qui sont ramifiés seulement aux \unsure{pts} de $X-U$.}

% lines 21-29
Un procédé de sections transversales prises ds $V$ et de ``courbe générique'' nous ramène au cas où $X$ est de dim. 1. \deleted{Alors le calcul explicite \unsure{revient} \unsure{th.}} \inserted{De plus} \deleted{\unsure{soit}} supposons que le modèle projectif régulier $\hat X$ de $X$ est lisse. La dualité nous ramène au cas de $\unsure{H^1_!|H^1_c}(X_{\overline K},\mathbb{Q}_\ell)$, la suite exacte au cas de $H^1_!(\hat X,\mathbb{Q}_\ell)=H^1(\hat X,\mathbb{Q}_\ell)$. On est donc ramené au cas d'une \uline{courbe} \inserted{\uline{régulière}} \uline{lisse} sur $K$. À expliciter en

\uline{Proposition.} 9.8. \unsure{. . . . . .}
% archive sheet number: 50
\scanpage{53}
% heading and introduction
\uline{Exposé II} \inserted{Cas d'action \deleted{\uline{\unsure{négligeable}}} du groupe fondamental} (Une variante du th. de monodromie) \deleted{\unsure{pour}}

Nous donnons dans le \deleted{\unsure{les articles|les autres}} présent exposé un résultat de nature analogue au théorème de monodromie de l'exposé précédent, \deleted{\unsure{pour}} à titre de nouvelle illustration de la méthode qui y était employée. Il ne servira plus dans la suite de ce séminaire.

\uline{1. Énoncé du résultat, et Corollaires.}

\deleted{\unsure{Énon}} Le résultat \deleted{essentiel} \inserted{principal du présent exposé} est le suivant.

% theorem, lines 7-12
\uline{Théorème} 1.1. Soient $k$ un corps \deleted{\unsure{transcendant} \unsure{quelconque}, tel que le \unsure{groupe} \unsure{soit} \unsure{de} \unsure{caract.} de \unsure{$k$}} \deleted{\unsure{sous-extension de type fini algébriquement}} de type fini (au sens absolu), $X$ un schéma sur \inserted{$S=\operatorname{Spec}k$} $k$, normal, localement de type fini, \deleted{\unsure{et}} géométriquement connexe, $\unsure{\xi|\zeta}$ un point géométrique \inserted{de $X$}, d'où (SGA 1 \uline{VIII}) une suite exacte
\[
1\longrightarrow\pi_1(\overline X,\xi)\longrightarrow\pi_1(X,\xi)\longrightarrow\pi_1(S,\xi)\longrightarrow1
\tag*{(1.1.1)}
\]
% lines 14-18
avec \deleted{\unsure{$\overline k$}} $\overline X=X\otimes_k\overline k$, $\overline k$ étant la clôture séparable de $k$ dans $k(\xi)$, (et $\xi$ désignant aussi les points géométriques de $S$ resp. $\overline X$ déduits canoniquement du \deleted{\unsure{$\xi$}} point géométrique $\xi$ de $X$). \deleted{Soit \unsure{$I$} \unsure{le} \unsure{groupe} \unsure{Considérons le plus grand}}

% lines 19-23
\deleted{\unsure{que}} Soient $\ell$ un nombre premier distinct de la caractéristique de $k$, \deleted{\unsure{et}} $\pi$ un sous-groupe ouvert de $\pi_1(X,\xi)$, $\pi_0$ son image dans $\pi_1(S,\xi)$, $\overline I=\pi\cap\pi_1(\overline X,\xi)$, de sorte qu'on a une suite exacte
\[
1\mathrel{\text{---}}\overline I\mathrel{\text{---}}\pi\mathrel{\text{---}}\pi_0\longrightarrow1.
\tag*{(1.1.2)}
\]
% lines 24-27
Soit $I_{\deleted{\unsure{\ell}}}(\ell)$ la composante $\ell$-primaire du plus grand quotient profini abélien $I_{\mathrm{ab}}$ de $I$, \deleted{Alors} d'où \inserted{par passage au quotient} une extension
\[
1\mathrel{\text{---}}I(\ell)\mathrel{\text{---}}\pi'_{\deleted{\unsure{\ell}}}\longrightarrow\pi_0\mathrel{\text{---}}1.
\tag*{(1.1.3)}
\]
% lines 28-32
Alors \deleted{$\unsure{[\pi',I(\ell)]}$ \unsure{est un groupe}} $I(\ell)$ est un module de type fini sur $\mathbb{Z}_\ell$, et $\overline{[\pi',I(\ell)]}$ \deleted{est de} d'indice fini dans $I(\ell)$, en d'autres termes \deleted{\unsure{$I(\ell)$}} le plus grand quotient profini de $I(\ell)$ sur lequel $\pi_0$ opère trivialement est un groupe \uline{fini}.
% archive sheet number: 52
\scanpage{54}
\mspage{II 2}
% lines 1-6
Ce théorème sera démontré dans les sections suivantes, \deleted{par réduction au cas où $X$ est un ouvert d'une courbe algébrique projective et lisse sur $k$, auquel cas l'énoncé \unsure{résultera|résulte} \inserted{\unsure{du résultat classique déduit}} de l'hyp. de Riemann pour les courbes algébriques sur un corps fini, démontrée par A. Weil \unsure{que} dans ce cas.}

% lines 7-10
Nous allons dans le présent n° donner quelques conséquences. Notons que la suite exacte (1.1) implique une suite exacte
\[
(I^{\mathrm{ab}}_{\deleted{\unsure{\mathrm{ab}}}})_{\pi_0}\longrightarrow\pi^{\mathrm{ab}}_{\deleted{\unsure{\mathrm{ab}}}}\longrightarrow\pi_0^{\mathrm{ab}}\longrightarrow1
\]
\deleted{\unsure{et}} et en passant aux composantes $\ell$-primaires
\[
I(\ell)_{\pi_0}\longrightarrow\pi^{\mathrm{ab}}(\ell)\mathrel{\text{---}}\pi_0^{\mathrm{ab}}(\ell)\mathrel{\text{---}}1,
\]
d'où on déduit :

% lines 14-17
\uline{Corollaire} 1.2. L'homomorphisme \inserted{surjectif} canonique
\[
\pi^{\mathrm{ab}}(\ell)\longrightarrow\pi_0^{\mathrm{ab}}(\ell)
\]
\deleted{\unsure{est}} a un \uline{noyau fini}.

\deleted{Utilisant I 6.1, on en déduit de 1.1 le}

% lines 19-27
\uline{Corollaire} 1.3. Soit
\[
u:\pi_1(X,\unsure{\xi|\zeta})\longrightarrow\operatorname{Aut}_{\mathbb{Q}_\ell}(\unsure{E|\mathbb{E}}\unsure{]|)}
\]
une représentation $\ell$-adique \deleted{\unsure{de $\pi$}}, tel que la représentation induite sur \deleted{\unsure{$I$}} \inserted{$\pi_1(\overline X,\xi)$} soit \inserted{\uline{essentiellement résoluble}} \deleted{\uline{abélienne}, \unsure{ou plus} généralement \deleted{\uline{commutative}} \uline{\unsure{abélienne}}} (i.e. $\exists$ un sous-groupe ouvert de $I$ dont l'action soit \deleted{abélienne} \inserted{résoluble}). Alors la représentation induite sur $I$ est \uline{essentiellement unipotente}, \deleted{et que} (donc essentiellement triviale si elle est semi-simple).
% archive sheet number: 53
\scanpage{55}
\mspage{II 3}
% lines 1-9
En effet, \deleted{\unsure{cette} \unsure{représentation} \unsure{passe à un sous-groupe assez petit}} \inserted{l'hypothèse \inserted{\unsure{indique|implique} que} il existe un ss-groupe} ouvert $\pi$ de $\pi_1(X,\unsure{\xi|\zeta})$ tel que la représentation induite sur $I_{\unsure{\pi|1}}=\pi\cap\pi_1(\overline X,\xi)$ soit abélienne. Prenant $\unsure{U|\pi}$ assez petit, on peut supposer (I \quad) que \deleted{elle se \unsure{factorise par}} l'image de \deleted{\unsure{$\pi$}} dans $\operatorname{Aut}_{\mathbb{Q}_\ell}(\unsure{E|\mathbb{E}})$ est un pro-$\ell$-groupe. Donc la représentation \inserted{de $\pi$} provient d'une représentation de $\pi'$, \deleted{\unsure{extension}} l'extension (1.3) de $\pi_0$ par $I(\ell)$. On conclut par I 6.1.

% lines 10-18
\uline{Corollaire} 1.4. Avec les notations de 1.3, supposons l'action de $\pi_1(X,\xi)$ lui-même \deleted{\unsure{semi}} \inserted{essentiellement} abélienne. Alors \inserted{la représentation induite sur $I$ \deleted{\unsure{alors}} est essentiellement triviale,} \deleted{\unsure{En particulier}} \inserted{i.e.} il existe un \deleted{\unsure{revêtement galoisien} fini surjectif} \inserted{sous-groupe \deleted{distingué} ouvert $U$} \deleted{\unsure{connexe} $X'\mathrel{\text{---}}X$,} de $\pi$ tel que la restriction de la représentation envisagée à $\unsure{\pi|U}$ \deleted{\unsure{par}} se factorise en une représentation de \deleted{\unsure{$K$}} $\pi_0=\operatorname{Im}(\unsure{\pi_q|\pi}\mathrel{\text{---}}\pi_1(S,\xi))$.

% lines 19-29
\deleted{Cela résulte du fait que :}

On prend \deleted{\unsure{$U$}} d'abord $\pi$ tel que la restriction à $\pi$ soit abélienne, \deleted{donc par} et que \deleted{\unsure{provenant de}} l'image de $\pi$ dans $\operatorname{Aut}_{\mathbb{Q}_\ell}(E)$ soit un pro-$\ell$-groupe. Alors la représentation provient d'une représentation de $\pi^{\mathrm{ab}}(\ell)$, et \deleted{il \inserted{donc} suffit d'appliquer} 1.2 nous assure qu'il existe un sous-groupe ouvert $U$ de $\pi^{\mathrm{ab}}(\ell)$ tel que $U\mathrel{\text{---}}\pi_0^{\mathrm{ab}}(\ell)$ soit injectif, donc la représentation restreinte à $U$ provient d'une représentation de son image dans $\pi_0^{\mathrm{ab}}(\ell)$. Il suffit
% archive sheet number: 54
\scanpage{56}
\mspage{II. 4}
% lines 1-4
alors de \deleted{considérer} \inserted{prendre} l'image inverse $V$ de $U$ dans $\pi$, la restriction de la \unsure{$u$} à $V$ provient d'une représentation de l'image de $V$ dans $\pi_0^{\mathrm{ab}}(\ell)$, donc d'une repr. de l'image de $V$ dans $\pi_0$, cqfd.

\diagram{\begin{tikzcd}[ampersand replacement=\&,column sep=small,row sep=small]
\& V\arrow[r,hook]\arrow[d]\arrow[dl] \& \pi\arrow[d]\arrow[dll]\\
\operatorname{Im}V\subset\pi_0\arrow[dr]\arrow[drr] \& U\arrow[r,hook]\arrow[d,"\unsure{\wr|\simeq}" description] \& \pi^{\mathrm{ab}}(\ell)\arrow[d]\\
\& U\arrow[r,hook] \& \pi_0^{\mathrm{ab}}(\ell)
\end{tikzcd}}

% lines 5-6
On va donner une application de ces résultats à la théorie des V.A.

% lines 7-12, rewritten statement
\uline{Théorème} 1.5. Soient $k$, \deleted{$X,\unsure{\xi|\zeta}$ comme dans 1.1, \unsure{et}} \inserted{un corps sép. clos, $X$ schéma loc. de type fini sur $k$, normal et \deleted{\unsure{géométriquement}} \inserted{géométriquement intègre}, $\xi$ est pt générique de $X$,} \deleted{\unsure{sous} \inserted{en supposant le corps de type fini mais sép. clos}} \deleted{\unsure{que}} $A$ un schéma abélien sur $X$, et \deleted{\unsure{soit}} $\ell$ un nombre premier distinct de la caractéristique de $k$. Supposons que la représentation de $\pi_1(\overline X,\xi)$ dans $T_\ell(A_\xi)$ soit \deleted{\uline{\unsure{semi-abélienne}}} \inserted{essentiellement résoluble.}

\diagram{\begin{tikzcd}[ampersand replacement=\&]
A\arrow[d,dash]\\X\arrow[d,dash]\\S
\end{tikzcd}}

% lines 13-21
Alors cette représentation est essentiellement \deleted{unipotente, donc essentiellement triviale si elle est semi-simple $(*)$. Dans ce \unsure{cas},} \inserted{triviale, et} il existe un revêtement étale fini surjectif $X'\longrightarrow X$, \deleted{et une extension finie \unsure{de type} finie séparable $k'$ de $k$, \unsure{dont} \unsure{un} $k$-morphisme $X'\longrightarrow\operatorname{Spec}k'$ \unsure{fais}} un schéma abélien $B$ sur $k$, \deleted{tel que $A_{X'}$} et une \uline{isogénie} \inserted{\uline{radicielle}} $B_{X'}\mathrel{\text{---}}A_{X'}$. \deleted{qui soit radicielle} En particulier, si $k$ est de car. nulle, \deleted{\unsure{donc}} $A_{X'}$ est isom. à $B_{X'}$ (i.e. est un schéma abélien constant).

% lines 22-25
Prouvons d'abord l'assertion que la représentation de $\pi_1(X,\xi)$ dans $T_\ell(A_\xi)$ est \deleted{unipotente \unsure{en car.}} \inserted{ess. unipotente}. \deleted{Dans} \deleted{\unsure{ceci}} Soit $U$ un voisinage ouvert affine de $\xi$,

% footnote, cancelled
\deleted{$(*)$ \unsure{Il y a tout lieu de conjecturer} que sans hypothèse supplémentaire l'action de $\pi_1(\overline X,\xi)$ sur $T_\ell(A_\xi)$ est \uline{toujours} \unsure{$\frac12$ simple|semi-simple}.}
% archive sheet number: 55
\scanpage{57}
\mspage{II. 5}
% lines 1-4
alors, \deleted{\unsure{$k$}} $U$ étant normal, $\pi_1(U,\unsure{\xi|\zeta})\mathrel{\text{---}}\pi_1(X,\xi)$ est surjectif, donc il \deleted{est ramené au} suffit de prouver l'assertion pour $U,\xi$ et $A|U$. Cela nous permet de supposer $X$ affine, donc de t.f. sur $k$.

% lines 5-9
On sait alors que $X,A$ provient de $S_0=\operatorname{Spec}k_0$, $k_0$ un sous-corps de t.f. de $k$, $X_0$ schéma de t.f. sur $S_0$, $A_0$ schéma abélien sur $X_0$. Alors, \deleted{\unsure{$\xi$ définit un pt}} $X\mathrel{\text{---}}X_0$ étant fid. plat et $X$ normal, $X_0$ est normal.

\diagram{\begin{tikzcd}[ampersand replacement=\&]
A_0\arrow[r,dash]\arrow[d,dash] \& A\arrow[d,dash]\\
X_0\arrow[r,dash]\arrow[d,dash] \& X\arrow[d,dash]\\
S_0\arrow[r,dash] \& S
\end{tikzcd}}

% lines 10-13
De plus, $\xi$ définit un pt géométrique de $X_0$, donc on peut considérer \deleted{l'hom. de} la suite exacte du type (1.1.1)
\[
1\mathrel{\text{---}}\pi_1(\overline X_0,\xi)\mathrel{\text{---}}\pi_1(X_0,\xi)\mathrel{\text{---}}\pi_{\unsure{0|1}}(S_0,\xi)\longrightarrow\unsure{0}.
\]
% lines 14-17
\deleted{De plus, on \unsure{désigne le}} $\overline k_0$ est la clôture séparable de $k_0$ dans $k$, on a $\overline X_0\simeq X_0\otimes_{k}\overline k_0$, donc $\unsure{X_0|X}\simeq\overline X_0\otimes_{\overline k_0}k$. On a un hom \deleted{\unsure{de}}
\[
\pi_1(X,\xi)\longrightarrow\pi_1(\overline X_0,\xi)
\tag*{(1.5.1)}
\]
qui est en tous cas surjectif \deleted{\unsure{$X_0$ et}} (réf \quad).

% lines 19-25
On va considérer la représentation naturelle de $\pi_1(X,\xi)$ dans $T_\ell(A_{0\xi})$ \inserted{$=T_\ell(A_\xi)$}, et la représentation induite via $\pi_1(X,\xi)\longrightarrow\pi_1(\overline X_0,\xi)\hookrightarrow\pi_1(X,\xi)$ n'est autre que celle qu'il s'agit d'étudier. \deleted{\unsure{À}} Par la surjectivité de (1.5.1), dire qu'elle \deleted{\unsure{est que}} \inserted{\unsure{est} ess. résoluble (resp.} \inserted{ess\textsuperscript{t}} unipotente) signifie que celle induite sur $\pi_1(\overline X_0,\xi)$ l'est. \deleted{Donc} On conclut donc grâce à 1.3.
% archive sheet number: 56
\scanpage{58}
\mspage{II. 6}
% top insertion, continued in slanted left margin
\inserted{Il existe donc un \unsure{ss-gro|ss-groupe} ouvert $\pi'\hookrightarrow\pi_1(X,\unsure{\xi|\zeta})$ sur lequel la repr. est unipotente. Mais $\pi'$ définit un rev. étale connexe $X'$ de $X$, \deleted{\unsure{tel}} $\xi'$ pt au-dessus tel que la représentation naturelle de $\pi_1(X',\xi')=\pi$ dans $T_\ell(A_{X'}(\xi'))$ soit unipotente. On peut donc, pour la démonstration, supposer déjà que $\pi_1(X,\xi)$ opère de façon unipotente.}
% slanted: continuation of the insertion in the left margin

% lines 3-8: crossed-out argument
\deleted{Supposons maintenant que la représentation envisagée de $\pi_1(X,\xi)$ soit aussi semi-simple, donc, étant unipotente, elle est triviale sur un ss-groupe \deleted{distingué} \inserted{ouvert} $U$ de $\pi_1(X,\xi)$. Ce sous-groupe définit un revêtement \inserted{connexe $\xi$-pointé} $X'$ de $X$, et par construction la représentation de $\pi_1(X',\xi)$ dans $T_\ell(A_{\xi'})$ est triviale.}

% lines 9-13
\deleted{\unsure{Cela}} Il suffit alors d'appliquer \inserted{l'adaptation \unsure{suivante|immédiate} d'}un résultat bien connu de \unsure{descente}, \deleted{\unsure{simple}} utilisant Mordell-Weil (cf Inventiones, Un th. sur les hom. de schémas abéliens, prop. 4.4) : \deleted{\unsure{on en déduit aussitôt que}}

\uline{Lemme} 1.5.1. \inserted{Si l'action de $\pi_1(X,\xi)$ sur $T_\ell(A(\xi))$ est \deleted{\unsure{ess.}} unipotente,} \deleted{\unsure{Si}} Alors cette action est triviale, et $\exists$ schéma abélien $B$ sur $k$, et \inserted{une isogénie radicielle $B_X\longrightarrow A$}.
% slanted: the isogeny insertion is in the left margin

% lines 14-18
\uline{Corollaire.} 1.6. \deleted{Sous les hyp.} Si pour \uline{un} $\ell\ne\operatorname{car}k$, l'action de $\pi_1(X,\xi)$ sur $T_\ell(A_\xi)$ est \deleted{unipotente \unsure{et}} \inserted{ess\textsuperscript{t} résoluble} \deleted{semi-simple} ($\Longleftrightarrow$ \deleted{\unsure{quasi}} essentiellement triviale) il en est de même pour tout autre $\ell'\ne\operatorname{car}k$.

% lines 19-26
\uline{Corollaire} 1.7. Supposons \deleted{que} \inserted{$\operatorname{car}k=0$, et} \deleted{\unsure{que}} l'action de $\pi_1(X,\xi)$ sur $T_\ell(A_\xi)$ soit \unsure{abélienne|unipotente} et semi-simple, et supposons qu'il existe un \inserted{entier $n\geq3$}, \deleted{\unsure{premier} $\ell'\ne\operatorname{car}k$,} tel que la représentation \deleted{de} induite \inserted{de} $\pi_1(X,\xi)$ sur ${}_n A(\xi)$ soit triviale, i.e. que ${}_n A$ soit un revêtement en groupes \uline{trivial} de $X$. Alors \deleted{\unsure{en 1.5 on peut prendre} $X'\longrightarrow X$, (en particulier si $\operatorname{car}k=0$)} $A$ est un schéma abélien \uline{constant} sur $X$.

% lines 27-29
En effet, \deleted{\unsure{alors} $\pi_1(X,\xi)$ opère sur $T_\ell(A_\xi)$ par l'intermédiaire d'un groupe \uline{fini}. Or on} la théorie de la descente galoisienne et la structure discrète de $\operatorname{Aut}_k(B)$ impliquent que
% archive sheet number: 57
\scanpage{59}
\mspage{II. 7}
% lines 1-6
$A$ est décrit par une représentation
\[
\Gamma\mathrel{\text{---}}\operatorname{Aut}(B)
\]
où $\Gamma$ est le groupe de Galois de $X'/X$. L'hypothèse \unsure{indique|implique} que $\Gamma$ opère trivialement sur ${}_n B(k)$, et on sait \unsure{q|que} (comme $n\geq3$) que cela \unsure{implique|entraîne} que l'action de $\Gamma$ sur $B$ est triviale, d'où la conclusion.

% lines 7-8
\uline{Remarque} 1.8. \deleted{Le théorème \unsure{précédent}} Par descente, on voit que 1.7 reste valable tel quel sans supposer $k$ \deleted{alg.} \inserted{sép.} clos.

% lines 9-12
L'\uline{énoncé} 1.5 est de nature purement géométrique, (bien que sa démonstration soit de nature arithmétique). Voici une variante arithmétique de 1.5 :

% lines 13-23
\uline{Théorème} 1.9. Soit $K$ un corps de type fini \deleted{\unsure{$A$}} $\overline K$ une clôture séparable de $K$, \inserted{$\pi=\operatorname{Gal}(\overline K/K)$,} $A_K$ un schéma abélien sur $K$. Supposons l'action de $\pi$ sur $T_\ell(A(\overline K))$ \deleted{abélienne} \inserted{essentiellement résoluble}. Alors, il existe une extension finie \deleted{\unsure{$B$}} \inserted{$k'\subset\overline K$} du corps premier, \deleted{et un} \inserted{et} un schéma abélien $B$ sur $k'$, et une isogénie radicielle \deleted{$B_{K'}\simeq A$}
\[
B\otimes_{k'}K'\simeq A\otimes_K K',
\]
et on a posé \deleted{\unsure{$\ell$}} $K'=Kk'$. Lorsque $\operatorname{car}K=0$, \deleted{\unsure{donc}} \inserted{l'action de $\pi$ \deleted{\unsure{sur}} sur $T_\ell(A(\overline K))$ est \uline{ess\textsuperscript{t} abélienne} $(*)$ ou tout au moins \unsure{ess\textsuperscript{t} nilpotente}} $k'$ est \deleted{de car nulle} un corps de nbs, on peut supposer que $B$ \deleted{\unsure{est} \unsure{$R'$} \unsure{fini}} provient d'un schéma abélien sur \inserted{$R$ \deleted{\unsure{$R'$}}} l'anneau des entiers de $k'$.

% footnote
$(*)$ Hypothèse sans doute \deleted{\unsure{conséquence}} \unsure{autom. satisfaite} \unsure{Donc à|Donc on} \unsure{démontrer que} l'action est tjs $\frac12$ simple.
% archive sheet number: 58
\scanpage{60}
\mspage{II. 8}
% heading and lines 1-6
1.10. \uline{\unsure{Démonstration.|Démonstr.}} \deleted{\unsure{Il suffit|En effet} \unsure{de|on voit de} 1.4.} \inserted{1°)} Existence de $B$.

Soit $k$ la \deleted{clôture} \inserted{fermeture} algébrique du corps premier dans $K$, \deleted{de sorte que $K$ est une extension régulière} \inserted{et soit $X$ un schéma \deleted{de type fini} \inserted{intègre} de type fini} \deleted{extension} sur $k$, \inserted{de corps des fractions $K$,} tel que $A_K$ provienne d'un schéma abélien sur $X$. \deleted{Alors $X$ est géom.} On peut supposer $X$ \deleted{géom.} normal (et $=$ régulier), et \deleted{il est} d'après la déf. de $k$, il est géom. intègre (NB. $k$ est parfait, étant une extension finie du corps premier).

\diagram{\begin{tikzcd}[ampersand replacement=\&]
X\arrow[r,dash]\arrow[d,dash] \& X'\arrow[d,dash]\\
k\arrow[r,dash] \& k'
\end{tikzcd}}

% lines 9-15
\deleted{\unsure{L'hyp.}} Alors l'extension $\overline K$ de $K$ définit un pt géom. $\unsure{\xi|\zeta}$ de $X$, et l'hyp. sur $A_K$ signifie aussi que la représentation de $\pi_1(X,\xi)$ dans $T_\ell(A_\xi)$ est ess\textsuperscript{t} \deleted{abélienne} \inserted{résoluble}. En vertu de \unsure{1.3|1.5}, \deleted{quitte à} passant à un rev. étale \inserted{surjectif} $X'$ de $X$ et à une extension finie $k'$ de $k$, on trouve que \inserted{$A_{X'}$ est \unsure{radiciellement} isogène à un $B_{k'}\otimes_{k'}X'$} \deleted{la représentation de $\pi_1(X',\xi)$ sur $T_\ell(A(\xi))$ est \unsure{définie} par l'intermédiaire d'une représentation de $\pi_1(\operatorname{Spec}(k'),\xi)$, \unsure{donc} i.e. $T_\ell(A_{X'})$ est un faisceau $\ell$-adique géom. trivial sur $X'$. En vertu du résultat déjà cité de \deleted{Inventiones} (\quad), il s'ensuit que l'on peut trouver $B$ sur $k'$ \deleted{tel que} et une isogénie radicielle \deleted{\unsure{$B\otimes$}} $B_{X'}\longrightarrow A_{X'}$,} ce qui achève la démonstration de l'hyp. en car. quelconque.

% lines 24-27
En car. $p>0$, comme $\operatorname{Gal}(\overline k'/k')\simeq\widehat{\mathbb{Z}}$ est abélien, on conclut que l'action de $\operatorname{Gal}(\overline K/K)$ sur $T_\ell(A(\overline K))$ est en fait \uline{abélienne}.

\deleted{\unsure{Démonstr.} \unsure{Lorsque la} caractéristique est nulle, \unsure{on en conclut que}}

% bottom marginal additions and next heading
\margin{Cas de l'action supposée ess\textsuperscript{t} abélienne}
% slanted
1.11. \uline{Bonne réduction} \inserted{aux pts de car. $\ne\ell$} Si $\overline k'$ est la clôture sép. de $k'$ dans $\overline K$, il est bien connu que les hom. can.
% archive sheet number: 59
\scanpage{61}
\mspage{II. 9}
\diagram{\begin{tikzcd}[ampersand replacement=\&]
\operatorname{Gal}(\unsure{\overline K'|\overline K}/K')\arrow[r]\arrow[d] \& \operatorname{Gal}(\overline k'/k')\\
\operatorname{Gal}(\overline K/K) \&
\end{tikzcd}}

% lines 4-16
ont des images qui sont des sous-groupes ouverts. Donc \deleted{l'hyp.} dire que $A_K$, ou $A_{K'}$, ou $B$ sur $k'$, donnent lieu à des actions ess\textsuperscript{t} \deleted{\unsure{résolubles}} \deleted{abéliennes} \inserted{abéliennes}, revient au même (Ceci est vrai chaque fois qu'on a une extension de type fini\unsure{\ldots}) ! \deleted{\unsure{Donc}} On en conclut p.ex. déjà une réciproque \deleted{\unsure{donc la}} de l'énoncé donné, lorsque $\operatorname{car} K>0$ \inserted{donc $k'$ est un corps fini} : s'il existe $B/k'$ avec $B\otimes_{k'}K'$ isogène à $A\otimes_K K'$, alors l'action de \deleted{$\pi(K/K')$ sur} $\operatorname{Gal}(\overline K/K)$ \deleted{sur} $T_\ell(A(\overline K))$ est \deleted{\unsure{alors}} ess\textsuperscript{t} abélienne : en effet, celle de \deleted{\unsure{Gal}} \deleted{\unsure{$\xi$}} $\operatorname{Gal}(\overline k'/k')$ \inserted{$\simeq\widehat{\mathbb{Z}}$} sur $T_\ell(B(\overline k'))$ l'est évidemment !

% lines 17-18
Dans le cas en car $K=0$ cela nous montre que

% lines 19-29
Pour prouver la dernière assertion de 1.9, on est ramené \deleted{\unsure{à pr.}} au cas où $K$ est déjà un corps de nombres. Dans ce cas, \deleted{d'après le} \deleted{\unsure{théorème}} \deleted{\unsure{critère de}} ce résultat est bien connu. Par \deleted{\unsure{le critère de Néron}} \deleted{\unsure{par}} le critère de Néron-\inserted{Ogg}-\unsure{Safarevitch|Schafarevitch} (\unsure{[\quad \S\ S., th.\ 2]|[\quad \S\ S., T.\ 2]}, cf aussi Exp.\quad) on est ramené à voir que si $R$ est l'anneau des entiers de $K$, alors \inserted{pour tt idéal maximal $\mathfrak m$ de $R$,} l'action du groupe d'inertie $I\subset\pi=\operatorname{Gal}(\overline K/K)$ relatif à $R_{\mathfrak m}$ est quasi-\deleted{unipotente}\inserted{triviale} (de sorte qu'après une extension finie convenable \inserted{$K'$} de $K$, on
\margin{\uline{développer}\ \uline{un peu}}
% slanted
% archive sheet number: 60
\scanpage{62}
\mspage{II. 10}
% lines 1-3
trouve une bonne réduction en les \deleted{\unsure{tous points de}} ``places'' de $K'$ au dessus de la place de $K$ définie par $R$). On a le

% lines 4-10
\uline{Lemme} 1.12. Soit $V$ un anneau de val. discrète, \inserted{de corps résiduel $k$,} \deleted{\unsure{dont le corps}} $\ell$ un nb premier $\ne\operatorname{car} k$, supposons que $k$ satisfasse l'hyp. de I 6.2. \inserted{les notations sont celles de I 1.} Alors pour tte représentation \inserted{$\ell$-adique} essentiellement abélienne de $\pi=\operatorname{Gal}(\overline K/K)$, la restriction à $I(\ell)$ est essentiellement triviale.

% lines 11-12
Se déduit de I 6.2 comme \deleted{\unsure{l}} on vient de déduire \unsure{1.4|1.3} de 1.1. via 1.2.

\margin{1) Pour le cas des pts de car. $\ell$, cf feuille \unsure{``le pb des puissances entières''|``le pb des poids entiers''} p. 13}

% lines 14-21
\uline{Remarque} 1.13. \deleted{\unsure{a}} Le seul exemple connu où on sache vérifier l'hypothèse \deleted{\unsure{de sem}} de l'action \deleted{abélienne} \inserted{\uline{abélienne et}\ \uline{semi-simple}} d'un groupe de Galois sur \inserted{un} $T_\ell(A_{\overline K})$ est celui où $A$ est une variété ``de type CM'' au sens de Shimura-Taniyama [\quad]. \inserted{Nous voyons donc que} \deleted{Dans} \unsure{ce cas} (le corps de définition de $A$ étant quelconque) $A$ devient radiciellement isogène, sur une extension \deleted{finie} séparable finie convenable, à \unsure{une|un} schéma abélien définissable sur un corps fini (cas car. $p>0$), resp. sur l'anneau des entiers d'un corps de nombres (cas de la car. $p=0$). Ce dernier cas \inserted{tout au moins} était bien connu [\quad].

% final line
Il semble plausible qu'en caractéristique
% archive sheet number: 61
\scanpage{63}
\mspage{II. 11}
% lines 1-4
nulle, \deleted{\unsure{un}} un schéma abélien sur un corps de type fini $K$ (qu'on peut $=$ supposer un corps de nbres) tel que l'action de $\operatorname{Gal}(\overline K/K)$ sur $T_\ell(A(\overline K))$ soit ess\textsuperscript{t} abélienne, \inserted{(ou seulement ess. résoluble)} soit de type CM.

% lines 5-12
1.14. On notera cependant que \deleted{les} \inserted{certains des} arguments donnés dans le présent numéro dépassent le cadre des schémas abéliens, et s'appliquent \deleted{\unsure{aussi bien}} \inserted{p.ex.} à des ``motifs'' de \deleted{degré} \inserted{poids} quelconque (une fois la théorie des motifs assise sur des bases solides). Notons seulement qu'une démonstration \unsure{de réduction|par réduction} à une situation arithmétique, ess. identique à la démonstration de 1.5, permet de prouver ceci :

% lines 13-22 and corrections between lines
\uline{Théorème} 1.15. Soient $k$ un corps sép. clos, $Y$ un schéma \deleted{connexe} \deleted{de t.f. sur $k$, \unsure{normal}} loc de t.f. sur $k$, \deleted{\unsure{qui est}} connexe et normal, $\unsure{\xi|\zeta}$ un pt géométrique de $Y$, $f:X\longrightarrow Y$ un morphisme propre \deleted{\unsure{et lisse}}, $\ell$ un nb premier $\ne\operatorname{car} k$, \deleted{\unsure{supposons}} $i$ un entier, \inserted{supposons que \deleted{le faisceau} $R^i f_*(\mathbb{Z}_\ell)$ soit un faisceau $\ell$-adique \unsure{const\textsuperscript{t}|const.} sur $Y$)} Considérons la représentation de $\pi_1(X,\xi)$ dans $E=H^i(X_\xi,\mathbb{Q}_\ell)$, \inserted{et le groupe algébrique $G\subset\operatorname{Aut}_{\mathbb{Q}_\ell}(E)$ engendré. Alors le} \deleted{Si cette représentation est \unsure{abél.} ess\unsure{.}} \deleted{\unsure{est}} \inserted{radical $R$ de $G^\circ$ est unipotent. En particulier, si} \deleted{abélienne, elle est ess. \unsure{unipotente}} \inserted{la représentation envisagée de $\pi_1(X,\xi)$ est ess. résoluble, elle est ess. unipotente; de + si on la suppose de plus semi-simple $^{(*)}$, elle} \deleted{\unsure{elle est triviale}} est essentiellement triviale. \deleted{\unsure{elle l'est}}
\margin{\unsure{tordu|tordus} (\unsure{tjs|tjrs} vrai à condition \deleted{de} \unsure{en réduisant|réduire|un réduit} $Y$)}
% slanted; note directed at the inserted hypothesis before Considérons

% lines 23-26
1.16. Dans l'énoncé 1.15, l'hypothèse de projectivité ne nous a servi que pour assurer que \deleted{\unsure{le foncteur de}} $R^i f_*(\mathbb{Z}_{\ell X})$ est un faisceau $\ell$-adique \uline{constructible} en voisinage du \unsure{point} générique $\unsure{\eta|y}$ de $Y$,

% footnote below rule
$(*)$ Ici \deleted{comme} comme dans 1.5, on conjecture \inserted{si $f$ est lisse} que l'hypothèse de $\frac12$ simplicité est automatiquement vérifiée.
% archive sheet number: 62
\scanpage{64}
\mspage{II. 12}
% lines 1-2, cancelled
\deleted{\unsure{constant tordu} sur $X$, dont la formation commute à \unsure{tout changement de la partie connexe à zéro}}

% lines 3-10
et que l'énoncé analogue est vrai après réduction de la situation à un corps de base de type fini. Cette hypothèse est \inserted{donc} inutile si on suppose qu'on dispose de la résolution des singularités pour les schémas de type fini \inserted{de dim $\leq\dim X_{\unsure{\overline\eta|\overline y}}$,} sur une clôture algébrique de $K$ ($\eta$ étant le pt générique de $Y$), \deleted{donc} \deleted{\unsure{elle est}} comparez I 8.8, \unsure{8.9 b}. Il en est en particulier ainsi si $\operatorname{car} k=0$, ou si $\dim X_\eta\leq2$.

% lines 11-16
On peut montrer que la conclusion est également valable si $i=1$, sans restriction \deleted{sur} de propreté sur $f$, \deleted{\unsure{dans ce cas cf}} en utilisant les réductions faites dans les n\textsuperscript{os} suivants; comparez n\textsuperscript{o} 5.

% lines 17-31, the whole remark crossed by diagonal strokes
\deleted{\uline{Remarque} 1.17. \deleted{Il est plausible que la forme} Avec les notations de 1.1, \deleted{\unsure{étant}} \deleted{\unsure{généralisée}} et soit $H$ un groupe \inserted{profini} quotient de $I$, invariant par l'action de $\pi$ par automorphismes intérieurs, et tel que $H$ soit un \deleted{\unsure{groupe analytique}} \deleted{\unsure{$\mathbb{Q}_\ell$}} \deleted{$\ell$-adique résoluble} pro-$\ell$-groupe \uuline{résoluble}. On trouve donc encore une extension}
\[
\deleted{1\mathrel{\text{---}}H\mathrel{\text{---}}\pi'\longrightarrow\pi_0\mathrel{\text{---}}1}.
\]
\deleted{Ceci posé, \uline{il est} plausible que l'on ait encore $\overline{[\pi_0,H]}$ ouvert dans $\pi'$. Cela semblerait devoir résulter de l'hyp. de Riemann-Weil en dim. quelconque, comme 1.1. résulte de l'hyp. de R.W. en dim 1. S'il en était ainsi, on}
% archive sheet number: 63
\scanpage{65}
\mspage{II. 13}
\uline{Fin de la démonstration de} 1.9.

\uline{Bonne réduction} \deleted{\unsure{aux pts de car. $\ell$}} (alias cas \unsure{nilpotent|unipotent}).

% lines 3-4
a) L'action est semi-simple et abélienne.

On utilise les faits suivants :

% lines 5-7
\uline{Lemme} Si $A$ est un schéma abélien sur un corps fini $k$, l'action du frobénius \inserted{géométrique} $f$ sur $T_\ell(A(\overline k))$ est semi-simple.

% lines 8-15
En effet, \deleted{\unsure{or}} \deleted{\unsure{c'est un}} \inserted{cette action est \deleted{précisément} \inserted{inverse de} celle du frobénius géométrique} \uline{endomorphisme} de $A_{\unsure{\overline k|k|\xi}}$. \deleted{Donc} Or cet $f$ appartient au \uline{centre} de l'anneau $\operatorname{End}(A)\otimes_{\mathbb{Z}}\mathbb{Q}$ des endomorphismes, \inserted{tensorisé par $\mathbb{Q}$,} qui, par Weil [2] est une algèbre \inserted{finie} semi-simple sur \deleted{\unsure{A}} $\mathbb{Q}$, donc son centre $Z$ est une \deleted{\unsure{extension}} algèbre finie \deleted{\unsure{séparable}} \inserted{étale} de $\mathbb{Q}$. \inserted{Donc $Z_\ell=Z\otimes_{\mathbb{Q}}\mathbb{Q}_\ell$ \deleted{\unsure{opère}} est une algèbre finie étale sur $\mathbb{Q}_\ell$;} \deleted{Donc} \deleted{$Z$ opère} \unsure{alors} il est clair $Z$ opère de façon semi-simple sur $T_\ell(A)\otimes_{\mathbb{Z}_\ell}\mathbb{Q}_\ell$,

% lines 16-24
\deleted{\uline{Lemme}} \inserted{\uline{Corollaire} 1} Soit donc $S$ un schéma \inserted{loc} de t.f. sur $\operatorname{Spec}\mathbb{Z}$, $\unsure{\xi|\zeta}$ un pt géom. de $S$, $A$ un schéma abélien sur $S$, \inserted{$\ell$ un nb premier distinct des car. résiduelles de $S$,} d'où une représentation de $\pi_1(S,\xi)$ sur \inserted{$E=$} $T_\ell(A(\xi))\otimes_{\mathbb{Z}_\ell}\mathbb{Q}_\ell$.

Pour tout point fermé $s\in S$, les éléments de frobénius $f_s\in\pi_1(S,\xi)$ [déterminé à \deleted{\unsure{auto}} conj. près comme l'image du gén. de frob. de $\pi_1(s)\deleted{\unsure{\xi}}$ par $\pi_1(s)\mathrel{\text{---}}\pi_1(S)$] opère \deleted{sur} de façon semi-simple.

% lines 25-27
\uline{Corollaire} 2 Si la représentation de $\pi_1(S,\xi)$ sur $E$ est \deleted{\unsure{si}} ess\textsuperscript{t} nilpotente, elle est \deleted{abélienne et} semi-simple et ess. abélienne.
% archive sheet number: 64
\scanpage{66}
\mspage{II. 14}
% lines 1-12
En effet, le groupe algébrique \inserted{$G\subset\operatorname{Aut}_{\mathbb{Q}_\ell}(\unsure{E_\ell|E})$} engendré par cette représentation est tel que $G^\circ$ est nilpotent, \deleted{Donc} \deleted{\unsure{par}} de \inserted{Donc :} un \uline{produit} \inserted{\uline{direct}} $T\times U$, $T$ un tore, $U$ unipotent. \inserted{\uline{A prouver} $U=1$.} Soit $\pi'\subset\pi$ l'image inverse de $G^\circ$; il correspond à un rev. \deleted{fini} étale connexe $S'$ de $S$, et quitte à passer à celui-ci, on peut supposer $G=G^\circ$. On utilise le th. de \uline{\unsure{Čebotareff|Čebotaref}} que \uline{les éléments de frobénius} \uline{sont denses} \uline{dans} $\pi_1(S,\unsure{\xi|\zeta})$. On en conclut que leurs images $f_s$ dans $G$ \uline{engendrent} $G$ comme groupe algébrique. Or $T$ est \unsure{l'ens} des élts semi-simples de $G$, donc $G=T$, i.e. $U=(1)$, cqfd.

% lines 13-19
b) \uline{Fin de la démonstration}. Il reste à prouver que \deleted{\unsure{pour tt}} pour tt $\mathfrak m\in\operatorname{Max}(R)$ tel que $R/\mathfrak m$ de car. $\ell$, $\exists$ extension finie $K'$ de $K$ telle que $A_{\deleted{\unsure{K}}K'}$ ait bonne réduction en les idéaux maximaux de $V'$ (où $V'$ est le normalisé de $V=R_{\mathfrak m}$ dans $K'$). Pour ceci, \inserted{on note que \deleted{l'action} en vertu \unsure{de|du} cor. 2, l'action de $\unsure{\operatorname{Gal}|\pi\mathrel{\text{.}}\operatorname{Gal}}(\overline K/K)$ sur $T_\ell(A(\overline K))\otimes_{\mathbb{Z}_\ell}\mathbb{Q}_\ell$ est $\frac12$ simple abélienne, donc il en est de \unsure{même}; \unsure{a fortiori} de celle du gr\textsuperscript{t} d'inertie $I_{\mathfrak m}$.}

% cancelled paragraph, lower block crossed diagonally
\deleted{\deleted{utilisant} un résultat d'un exposé ultérieur, on utilise le fait que \deleted{\unsure{il existe}} l'on peut trouver une extension \deleted{de} finie $K'/K$ tel que \deleted{le schéma} \inserted{le modèle de Néron de} $A_{K'}$ sur $R'$ ait ``réduction $\frac12$ st.R'' en tout point de \deleted{\unsure{$V'$}} $\operatorname{Spec}V'$. \deleted{Alors} \unsure{bien plus,}}
% archive sheet number: 65
